\section{Even strong traces and unweighted first roots}
\label{sec:even-strong-roots}
Two new fixed ordinary polynomials have fourteen unrestricted integer unknowns:
\[
 (\deg,\mathrm{ops})=(85,110),\qquad (87,108).
\]
The charged binary-operation counts are respectively $(63,24,23)$ and
$(62,23,23)$ for multiplication, addition and subtraction. They strictly improve
the matched $(86,112)$ and $(88,110)$ points; $(82,113)$ remains a separate tradeoff.
All six full fixed numeral payloads and their finite reconstruction recipe are
retained. These are ordinary reductions through the inherited native/compiler/
refutation theorem, not a consistency assertion, global optimum, or new
end-to-end proof-assistant certificate.

\subsection{Complete polynomial definitions}
The unknowns are
\[
 j,C,z,h,r,v,w,T,c,n,f,y,p_2,p_3\in\mathbb Z.
\]
Use the authentic ports $\beta,K_0,M_C,M_F,A_c,B_c$, with $M_F$ already shifted,
$\beta=2^{1064921}-1$, $A_c=2^{1113}$ and $B_c=2^{666104}$. Define
\begin{gather*}
 J=6j+1,\quad q_m=\beta J,\quad q=q_m+1,\quad W=12w+4,\quad X=qW,\quad Y=v,\\
 F=(K_0+W)C-zq_m,\quad
 R=[q(q-F)-C](q^2-1)+(M_C+qM_F)J,\\
 E=XY,\quad a=E+Y,\quad H=4a+3,\quad\Delta=a^2+H,\quad\ell=XY^2,\\
 \eta=c-Yk,\quad L=\ell k,\quad I=(8r+4)c^2,\quad Q=\Delta(f^2-1),\\
 V=f(cT-R)-c,\quad M=2q^3,\quad U=MH,\quad K=X(2X-3c)-2,\\
 Z_0=Un+MK+H(Y-q^3),\quad
 D_{\rm err}=(A_cF-q)^2+(B_cC-q)^2.
\end{gather*}
Choose the chart
\[
\begin{array}{c|c|c|c}
\deg&k&\mathcal B&Z\\\hline
85&R+1+h\ell&(Yk)^2(q^2-D_{\rm err})-q^2[4(Y\eta)^2+(kq^2)^2]&Z_0+UL\\
87&R+1+hE&(Yk)^2(q^2-D_{\rm err})-q^2[4(Y\eta)^2+(kq^3)^2]&Z_0
\end{array}
\]
and finish with
\begin{align*}
 G&=\mathcal B-1-p_2^2-p_3(p_3+1),&S&=f^2-\Delta I^2,\\
 u_u&=(V+y)^2,&E_a&=u_u-Q(u_u-V^2),\\
 D_\tau&=1+L(L+k),\\
 A&=U^2G,&B&=U^2D_\tau,\\
 N&=(Z^2-A-B)^2-4AB,&P&=(N+1)SE_a-1.
\end{align*}
The identity $E_a=V^2-(Q-1)y(2V+y)$ shows that sharing $u_u$ saves one
charged multiplication without changing the polynomial. Every name denotes a
polynomial abbreviation, not an additional unknown. The
87-chart evaluator uses $L=E(Yk)$ and does not separately compute $\ell$.

\subsection{Unweighted extraction, including all exceptional branches}
Unconditionally $q$ is even and nonzero, $|W|\ge4$, $|X|\ge8$, $H$ is odd and
nonzero, and $U\ne0$. Since $a=(X+1)Y$ is zero or has $|a|\ge7$,
$\Delta=(a+2)^2-1\ge3$. Also $\ell$ is even, so
$D_\tau=1+\ell(\ell+1)k^2\ge1$.

At $P=0$ the three integer factors are units. As $N$ is a square modulo four,
$N+1\ne-1$; as $4\mid I$, $S\equiv f^2\pmod {16}$ and $S\ne-1$.
Therefore $N=0$, $S=1$, $E_a=1$. With $\rho=Z/U$,
\[
 (\rho^2-G-D_\tau)^2=4GD_\tau.
\]
This forces $G\ge0$, including $G=0$. Hence $\mathcal B>0$ and $Yk\ne0$.
Since $|\ell|\ge8$ and $|X|\ge4|q|$,
\[
 D_\tau\ge1+\tfrac12X^2Y^4k^2
 \ge1+8q^2Y^4k^2>q^2Y^2k^2\ge\mathcal B>G.
\]
Thus $\rho\ne0$ before any division by it. The rationals
\[
 p=\frac{\rho^2+G-D_\tau}{2\rho},\qquad \tau=\rho-p
\]
satisfy $p^2=G$, $\tau^2=D_\tau$. Each is integral because a rational with
integral square is integral. This proof covers zero and coincident square classes;
no generic independence hypothesis occurs. The universal signed factorization is
\[
 (x^2-A-B)^2-4AB=\prod_{\epsilon_1,\epsilon_2\in\{\pm1\}}
 (x-\epsilon_1\sqrt A-\epsilon_2\sqrt B).
\]

Write $s=1$ for degree 85 and $s=0$ for degree 87. Restore the even-trace
fifteen-variable kernel by
\[
 p_1=p,\qquad n_{\rm old}=n+sL-(2L+k)-p+\tau.
\]
Its carrier is
\[
 \mathcal A=(Z_{\rm old}+Uk)^2+4U(LZ_{\rm old}-U)
           =(Z_{\rm old}+U(2L+k))^2-4U^2D_\tau,
\]
and therefore vanishes. The remaining rows are $S-1=0$, $E_a-1=0$ and
$\mathcal B-1-p_1^2-p_2^2-p_3(p_3+1)=0$. Conversely any zero of this kernel
has an integral $\tau$ with $Z_{\rm old}/U+2L+k=2\tau$: the rational on the
left has square $4D_\tau$, so is an even integer. Delete $p_1$ and set
\[
 n=n_{\rm old}-sL+(2L+k)+p_1-\tau.
\]
Then $\rho=p_1+\tau$ and $P=0$. These are total maps on zeros, not claimed
polynomial automorphisms; the root-dependent integral shift removes one gate.

\subsection{Signed unsquared trace: the new parity argument}
The comparison signed kernel has odd strong scale $(2r_{\rm old}+1)c^2$ and
squared trace $f_{\rm old}^2(cT_{\rm old}-R)-c$. Its retained proof is
\texttt{inherited\_round20/inherited/PROOF.md} for the $q^6/XY^2$ chart, and
\texttt{FREE\_SCALE15\_PROOF.md} in its inherited roots source for $q^8/XY$.
The normalized rank and auxiliary lemmas are in the same source's
\texttt{06\_PROOF.md}. We verify the changed hypotheses before restoring a
literal old zero.

The unchanged carrier and guard give $c\ne0$, the strict centered strip,
$q_0=|q|\ge16$, $x=|X|\ge4q_0$ and
\[
 q_0(q_0^2-1)<R<q_0^4-2.
\]
They give $|Y|>q_0^2$, $|XY^2|>3R+1$ in the 85 chart, and $|Y|>q_0^3$,
$|XY|>3R+1$ in the 87 chart. Carrier integrality gives even $Y$ and the exact
divisibilities $M\mid H(Y-q^3)$, $H\mid MK$ without premature cancellation.
Put $B_0=|a+2|$, $c_0=|c|$, $I_+=|8r+4|c_0^2$ and $f_+=|f|$.
Then $I_+>0$ is even, $B_0\ge3$, and
\[
 f_+^2-(B_0^2-1)I_+^2=1,\qquad f_+>2I_+.
\]
Pell classification gives $f_+=\chi_{B_0}(m)$, $I_+=\psi_{B_0}(m)$.
The residue $V\equiv-c\pmod {f_+}$ is unchanged. The ordinary odd auxiliary
index and folding lemma therefore recover $c_0=\psi_{B_0}(p)$ before any
residue modulo $c_0$ is needed: fold the auxiliary index to $0\le p\le m$,
and use $\psi_{B_0}(p)+c_0\le2I_+<f_+$ to turn its signed congruence into
equality. The first Pell classification and centered bounds give $p>n_0\ge24$
and $c_0>2p,2R$.

The normalized rank lemma gives $m=p\nu$ and $c_0\mid\nu$. Indeed, strong
divisibility gives $p\mid m$; with $D=\chi_{B_0}(p)$, composition gives
$I_+=c_0\psi_D(\nu)$, and
$\psi_D(\nu)\equiv \nu D^{\nu-1}\pmod {c_0}$ with $\gcd(D,c_0)=1$ gives $c_0\mid\nu$.
Now $\psi_{B_0}(m)\equiv m\pmod2$, so $m$ is even. If $c_0$ is even,
$\nu$ is even; if $c_0$ is odd, $p$ is odd and $m=p\nu$ again makes $\nu$ even.
As $D^2\equiv1\pmod {c_0}$ and $\chi_D(\nu)\equiv D^\nu\pmod {c_0}$,
\[
 |f|=\chi_D(\nu)\equiv1\pmod {c_0}.
\]
Consequently the unsquared trace has exactly the required signed residues
\[
 V\equiv-c\pmod {|f|},\qquad V\equiv\pm R\pmod {c_0}.
\]
The signed auxiliary step-down gives an odd index $\lambda$ with
$\lambda\equiv\pm p\pmod m$ and $\lambda\equiv\pm R\pmod {c_0}$.
Since $c_0\mid m$ and $0<p,R<c_0/2$, it follows that $p=R$.
Thus the extra sign does not invalidate the argument. The unchanged surviving
first-index branch gives $2n_0=R+1$. Whole-product projection, the modulo-three
orientation and late integer cancellation then give
\[
 c,k,Y>0,\quad X=2^R,\quad R\ \text{odd},\quad
 Y=(2v_{\rm scale}+1)q^3,\quad H\mid K,\quad0<\eta/k<1/4.
\]
These steps use only the already established indices, carrier and strip, not the
old square on $f$. In particular $B_0=a+2\equiv2\pmod8$ and $c$ is odd.

\subsection{Fresh strong witnesses: soundness and completeness}
One cannot keep $f,T$ when returning to an odd-scale source. Let
\[
 m_0=Rc,\qquad f_0=\chi_{B_0}(m_0),\qquad I_0=\psi_{B_0}(m_0).
\]
Composition gives $c^2\mid I_0$; as $R,c,m_0$ are odd, $I_0/c^2$ is odd.
Thus $r_{\rm old}=(I_0/c^2-1)/2$ restores the old strong row.
For every odd $R=2d+1$, the integer quotient polynomials obey
\[
 \chi_Z(2d+1)=ZQ_d(Z^2),\quad Q_d(0)=(-1)^d(2d+1),\quad
 Q_d(1-B_0^2)=(-1)^d\psi_{B_0}(2d+1).
\]
They follow from $Q_0=1$, $Q_1=4x-3$ and
$Q_{d+1}=(4x-2)Q_d-Q_{d-1}$. Set $S_0=\Delta I_0$ and
\[
 V_0=(-1)^{d+1}Q_d(S_0^2),\qquad u_0=\psi_{S_0}(R).
\]
Then $S_0^2V_0^2-(S_0^2-1)u_0^2=1$, while
$V_0\equiv-c\pmod {f_0^2}$ and $V_0\equiv-R\pmod c$.
Since $\gcd(c,f_0)=1$ and $f_0^2\equiv1\pmod c$,
\[
 T_{\rm old}=\frac{V_0+c+Rf_0^2}{cf_0^2},\qquad y_{\rm old}=u_0-V_0
\]
are integers and restore the squared auxiliary trace. Every outer, carrier and
guard coordinate is unchanged. This is a literal zero of the correct old signed
kernel (separately for each chart), not an assumed $R\equiv3\pmod4$ branch.
Its authenticated native/compiler theorem applies.

Conversely every old zero has the same established outer data. Double the seed:
\[
 f=2f_0^2-1,\qquad I=2f_0I_0.
\]
Because $B_0\equiv2\pmod4$ and $m_0$ is odd, $f_0\equiv2\pmod4$.
Hence $I/(4c^2)=(f_0/2)(I_0/c^2)$ is odd, so
$r=(I/(4c^2)-1)/2$ is integral and $I=(8r+4)c^2$.
The doubled pair has norm one; $f_0^2\equiv1\pmod c$ directly gives $f\equiv1\pmod c$.
Using the same signed quotient construction at $S=\Delta I$ yields $V_*,u_*$;
now
\[
 T=\frac{V_*+c+Rf}{cf},\qquad y=u_*-V_*
\]
are integers because the numerator is divisible by the coprime integers $c,f$.
This gives $V=f(cT-R)-c=V_*$ and the new auxiliary row. No positivity of $T$
is required on the unrestricted-integer domain. Apply the unweighted root map
to finish. Both directions are defined on all zeros using fresh witnesses.

\subsection{Attained degrees and exact coefficients}
Put
\[
 q_0=6\beta j,\ X_0=12q_0w,\ Y_0=v,\ a_0=X_0Y_0,\
 F_0=12wC-q_0z,\ R_0=-q_0^3F_0,\ \ell_0=X_0Y_0^2,\ U_0=8q_0^3a_0.
\]
Both complete strong/auxiliary tops are
\[
 [S]_{12}=-64a_0^2r^2c^4,\qquad [E_a]_{15}=2a_0^2f^3R_0y.
\]
For degree 87,
\[
 [N]_{60}=U_0^4R_0^4(q_0^8+\ell_0^2)^2,
\qquad
 [j^{56}w^8v^8z^5r^2c^4f^3y]P_{87}=-2^{91}3^{64}\beta^{56}\ne0.
\]
For degree 85 retain $k_0=R_0+h\ell_0$ and define
\[
 T_0=A_c^2Y_0^2F_0^2+4q_0^2Y_0^4+q_0^6.
\]
The entire cancellation
\[
 (Z_0+UL)^2-U^2D_\tau=Z_0^2+2ULZ_0-U^2(1+Lk)
\]
has degree 26. Thus the square in $N$ has degree 56 and the other term has top
\[
 [N]_{58}=4U_0^4\ell_0^2k_0^4T_0,
\qquad
 [j^{48}w^{10}v^{12}z^5r^2c^4f^3y]P_{85}=-2^{89}3^{58}\beta^{48}\ne0.
\]
These full leading forms, with respectively 18 and 70 monomials after expansion,
prove the exact degrees on the supplied fixed numeral slice, not merely generic
or sampled degrees.

\subsection{A separate correlated-root separation lemma}
Let $F,D\equiv1\pmod3$, $G\in\mathbb Z$, $3\mid Q$ and $b,c\in\mathbb Z$.
For algebraic roots
\[
 f^2=F,\quad\tau^2=D,\quad
 t^2=\frac{Q(fb-c)^2-1}{Q-1},\quad g^2=G,
\]
a rational $\rho=f+3\tau+9t+27g$ forces all four roots rational.
Work in the integral closure of $\mathbb Z_{(3)}$ in a normal field containing
all roots. They are integral over this ring, since $Q-1$ is a unit and
$t^2=\alpha+\beta f$ with $\alpha,\beta\in\mathbb Z_{(3)}$.
If an automorphism flips $f$, subtract the trace equations to obtain $f=3u$
with integral $u$. But $u^2=F/9$ is rational and integral over $\mathbb Z_{(3)}$,
hence lies in that ring, contradicting $F\equiv1\pmod3$. Thus all automorphisms
fix $f$, so $f$ is rational. Divide the conjugate difference by three to repeat
for $\tau$; then $t^2$ is rational and a three-adic unit, so divide by nine and
repeat for $t$. Finally $g$ is rational. This includes imaginary roots and
dependent square classes. A valid norm must use the different auxiliary roots
associated to each sign of $f$; the independent four-root norm cannot simply be
reused. The explored twelve-variable evaluators remained dominated, so no new
twelve-variable frontier claim is made.

\subsection{Verification and dependency boundary}
The standard-library checker binds every literal register independently, all
inputs and six ports, topology, coverage, liveness, exact operation counts,
four-factor identity, root maps, doubling, full leading cancellations and
attained coefficients. It rejects individual mutations of all 110 and 108 live
gates. Normal and optimized Python runs pass the same explicit conditions.
The finite reconstruction recomputes all 615883 coefficient bytes, 943 tiles and
69230 basis convolution checks. All 212 inherited source files are pinned.
Recurrence samples are diagnostics, not proofs of the unbounded lemmas above.
The current comparison snapshot is
\texttt{5274bc751a1566879ddc028a2d2a2af3c031654a}; original source attributions
and licenses are preserved. The ordinary native/compiler/refutation theorem
remains the named inherited dependency, not a claim inferred from finite checks.
